% =====================================================================
%  Chapter: Covariance, Invariance and Symmetry
%
%  Conventions used throughout:
%    * Gaussian units
%    * metric signature (+,-,-,-),  x^0 = ct,  epsilon^{0123} = +1
%    * Greek indices 0..3, Latin indices 1..3, summation convention
%
%  The file compiles as it stands (pdflatex, run twice).  To insert the
%  chapter into the main book, copy everything between the two markers
%  "CHAPTER BEGINS" / "CHAPTER ENDS" together with the macros and
%  theorem environments of the preamble.
% =====================================================================
\documentclass[11pt,a4paper,openany]{book}

\usepackage[T1]{fontenc}
\usepackage[utf8]{inputenc}
\usepackage{mathptmx}
\usepackage[margin=2.7cm]{geometry}
\usepackage{amsmath,amssymb,amsthm,mathtools,bm}
\usepackage{booktabs,array}
\usepackage{enumitem}
\usepackage{microtype}
\usepackage{tikz}
\usetikzlibrary{arrows.meta}
\usepackage[colorlinks=true,linkcolor=blue!50!black,citecolor=blue!50!black,urlcolor=blue!50!black]{hyperref}

% ---------------------------------------------------------------------
% theorem-like environments
% ---------------------------------------------------------------------
\theoremstyle{definition}
\newtheorem{definition}{Definition}[chapter]
\newtheorem{example}{Example}[chapter]
\theoremstyle{plain}
\newtheorem{proposition}{Proposition}[chapter]
\theoremstyle{remark}
\newtheorem{remark}{Remark}[chapter]

% ---------------------------------------------------------------------
% macros
% ---------------------------------------------------------------------
\DeclareRobustCommand{\hbar}{\hslash}
\newcommand{\dd}{\mathrm{d}}
\newcommand{\ii}{\mathrm{i}}
\newcommand{\pd}[2]{\frac{\partial #1}{\partial #2}}
\newcommand{\Lie}{\mathfrak{L}}
\newcommand{\eps}{\varepsilon}
\newcommand{\vx}{\bm{x}}
\newcommand{\vp}{\bm{p}}
\newcommand{\vE}{\bm{E}}
\newcommand{\vB}{\bm{B}}
\newcommand{\vL}{\bm{L}}
\newcommand{\vA}{\bm{A}}
\newcommand{\vn}{\bm{n}}
\newcommand{\vv}{\bm{v}}
\newcommand{\vbeta}{\bm{\beta}}
\newcommand{\Poi}[2]{\{#1,#2\}}
\newcommand{\weakeq}{\approx}
\DeclareMathOperator{\Stab}{Stab}
\DeclareMathOperator{\Diff}{Diff}
\DeclareMathOperator{\Isom}{Isom}
\DeclareMathOperator{\tr}{tr}

\begin{document}

% ===================== CHAPTER BEGINS ================================
\chapter{Covariance, Invariance and Symmetry}
\label{ch:cis}

\section{Three words, three levels of description}
\label{sec:three-words}

Few words in physics are used as loosely as \emph{invariant}, \emph{covariant}
and \emph{symmetric}. One reads of the ``Lorentz invariance of Maxwell's
equations'' and, a page later, of their ``Lorentz covariance''; of ``general
covariance'' as the foundation of general relativity and, in the next
paragraph, as a triviality of notation; of ``gauge symmetry'' that is plainly
not a symmetry in the sense in which rotations are. The three notions are
tied together by one logical structure, but they live at \emph{different
levels of description}, and mixing the levels is a dependable source of
confusion.

\begin{itemize}[leftmargin=1.6em,itemsep=2pt]
\item \textbf{Invariance} is a property of an \emph{object}: a number, a
  function, a tensor field, a functional, a particular solution. The object is
  left unchanged by a transformation.
\item \textbf{Covariance} is a property of a \emph{formulation}: a family of
  quantities, or an equation built from them, transforms in a definite,
  prescribed way (a representation of the transformation group), so that the
  equation keeps its form.
\item \textbf{Symmetry} is a property of a \emph{law} (a dynamical system): a
  transformation that carries solutions into solutions, or leaves the action
  unchanged.
\end{itemize}

Symmetry is the only one of the three that is a claim about the world. It
says that if an experiment on an isolated system is repeated after the
\emph{entire} arrangement has been displaced, rotated or boosted, the outcome
is the transformed outcome. That claim can be, and has been, tested with
extraordinary precision (Michelson--Morley-type experiments for boosts,
Hughes--Drever-type experiments for rotations). Covariance, by contrast, is
a statement about how we write equations; whether it carries any physical
content depends on what else is held fixed. The precise relation is the
central result of this chapter:
\begin{equation}\label{eq:slogan}
  \boxed{\;\text{symmetry}\;=\;\text{covariance of the equations}\;+\;
  \text{invariance of whatever is not dynamical}\;}
\end{equation}
Table~\ref{tab:dictionary} summarizes the vocabulary.

\begin{table}[ht]
\centering
\small
\renewcommand{\arraystretch}{1.3}
\setlength{\tabcolsep}{4pt}
\begin{tabular}{@{}p{2.1cm}p{3.9cm}p{4.6cm}p{4.0cm}@{}}
\toprule
Term & Property of & Formal statement & Question to ask\\
\midrule
Invariance & an object (scalar, tensor, solution, functional)
  & $g\cdot T = T$ & Is $T'$ equal to $T$?\\
Covariance & a formulation (a family of quantities, an equation)
  & $g\cdot T = D(g)\,T$ for a representation $D$
  & Does $E=0$ keep its form?\\
Symmetry & a law (dynamical system)
  & $g$ maps solutions to solutions, or $S[g\cdot\phi]=S[\phi]$
  & Is the transformed motion allowed?\\
\bottomrule
\end{tabular}
\caption{The three notions at a glance.}
\label{tab:dictionary}
\end{table}

The plan of the chapter is as follows. Section~\ref{sec:transformations}
fixes conventions and discusses what a transformation acts on.
Section~\ref{sec:definitions} gives the definitions and proves the relation
\eqref{eq:slogan}. Sections~\ref{sec:infinitesimal}--\ref{sec:noether}
develop the infinitesimal form, the distinction between the symmetry of a
\emph{law} and of a \emph{solution}, and Noether's two theorems.
Sections~\ref{sec:galilei}--\ref{sec:quantum} work through specific examples:
Newtonian mechanics, electrodynamics and the relativistic particle, general
relativity, gauge theories, the Kepler problem, and quantum mechanics;
Section~\ref{sec:kretschmann}, placed after general relativity, analyzes
Kretschmann's observation in detail.
Section~\ref{sec:summary} collects the results and lists the most common
misunderstandings; problems follow.

% ---------------------------------------------------------------------
\section{Transformations and what they act on}
\label{sec:transformations}

\subsection{Conventions}

We use Gaussian units, the metric $\eta_{\mu\nu}=\mathrm{diag}(+1,-1,-1,-1)$,
$x^0=ct$, and $\epsilon^{0123}=+1$. The Poincar\'e group acts as
\begin{equation}\label{eq:poincare}
  x'^\mu = \Lambda^\mu{}_{\nu}\,x^\nu + a^\mu,\qquad
  \Lambda^\mu{}_{\alpha}\Lambda^\nu{}_{\beta}\,\eta_{\mu\nu}=\eta_{\alpha\beta}.
\end{equation}
The electromagnetic potential is $A^\mu=(\phi,\vA)$ and the current is
$j^\mu=(c\rho,\bm j)$.

\subsection{Passive and active transformations}

A \emph{passive} transformation relabels the points of spacetime,
$x\mapsto x'=f(x)$, and asks how the \emph{same} physical field is described
in the new coordinates. For a scalar and for a vector field,
\begin{equation}\label{eq:passive}
  \phi'(x') = \phi(x),\qquad
  A'^\mu(x') = \frac{\partial x'^\mu}{\partial x^\nu}\,A^\nu(x).
\end{equation}
An \emph{active} transformation by $g$ moves the field itself, keeping the
coordinates: for $g\colon x\mapsto \Lambda x + a$,
\begin{equation}\label{eq:active}
  (g\cdot\phi)(x) = \phi(g^{-1}x),\qquad
  (g\cdot A)^\mu(x) = \Lambda^\mu{}_\nu\,A^\nu(g^{-1}x).
\end{equation}
The new field \emph{function} is given by the same formula in both cases
(solve \eqref{eq:passive} for $\phi'(y)=\phi(f^{-1}y)$ and compare with
\eqref{eq:active}); what differs is the interpretation. In the passive
reading one compares two descriptions of one physical situation; in the
active reading one compares two different physical situations in one
description. Symmetry is an active notion: it asks whether the second
situation is also allowed by the law. Covariance is naturally passive: it asks
whether the \emph{description} keeps its form.

\subsection{What must be transformed}

The active reading carries a warning that will turn out to be the heart of
the matter: \emph{everything that enters the equations must be transformed,
or else be invariant.}

\begin{example}[Coulomb field of a fixed charge]\label{ex:coulomb}
The equation $\nabla\!\cdot\vE=4\pi\rho$ with a prescribed source
$\rho=q\,\delta^3(\vx-\vx_0)$ has the solution
$\vE(\vx)=q(\vx-\vx_0)/|\vx-\vx_0|^3$. Translating the field alone,
$\vE_a(\vx)=\vE(\vx-\bm a)$, gives a field whose divergence is
$4\pi q\,\delta^3(\vx-\vx_0-\bm a)$: it is \emph{not} a solution for the
original source. If the source is translated too, the pair
$(\vE_a,\rho_a)$ is a solution. Translation is therefore a symmetry of the
closed system ``field $+$ charge,'' and is broken as soon as the charge is
treated as a fixed, external, non-dynamical object.
\end{example}

Whatever is held fixed and not transformed is called a \emph{background} (or
\emph{absolute}) structure: a prescribed source, a fixed metric, a preferred
velocity field. Prescribing backgrounds is unavoidable in practice, and it is
precisely what separates covariance from symmetry.

\subsection{Representations}

A group $G$ acts on a vector space $V$ through a \emph{representation},
a map $D\colon G\to \mathrm{GL}(V)$ with $D(g_1g_2)=D(g_1)D(g_2)$.
Fields are classified by the representation of the spacetime group under
which they transform: scalars (trivial), vectors, tensors, spinors. A scalar
field transforms by the first formula of \eqref{eq:active}, a vector field by
the second, with $D(g)=\Lambda$. The trivial representation $D(g)=1$ is
special: objects transforming in it are \emph{invariant}. Invariance is thus
the limiting case of covariance in which the representation is trivial.

% ---------------------------------------------------------------------
\section{Definitions and the basic theorem}
\label{sec:definitions}

Let $\phi$ denote the dynamical fields (collectively) and $b$ the background
fields, both defined on spacetime $M$ (for a mechanical system $M$ is the time axis;
see Section~\ref{sec:dof}). A group $\mathcal G$ acts on pairs,
$g\colon(\phi,b)\mapsto(g\cdot\phi,\,g\cdot b)$. A system of field equations
is an operator $E$,
\begin{equation}\label{eq:EOM}
  E[\phi;\,b]=0 .
\end{equation}

\begin{definition}[Invariance]\label{def:invariance}
An object $T$ (a field, a functional, a tensor, a solution) is
\emph{invariant} under a group $G$ if $g\cdot T=T$ for all $g\in G$.
\end{definition}

Familiar examples: the metric $\eta_{\mu\nu}$ is invariant under the
Lorentz group, which is exactly the content of \eqref{eq:poincare}; $\delta_{ij}$
and $\epsilon_{ijk}$ are invariant under $SO(3)$; the interval
$\eta_{\mu\nu}\Delta x^\mu\Delta x^\nu$ and the mass-shell value
$p_\mu p^\mu=m^2c^2$ are Lorentz invariants.

\begin{definition}[Covariance]\label{def:covariance}
A family of quantities $T^A$ is \emph{covariant} (transforms covariantly)
under $G$ if
\begin{equation}\label{eq:covariant-object}
  g\cdot T^A = D(g)^A{}_B\,T^B
\end{equation}
for a representation $D$ of $G$. The equations \eqref{eq:EOM} are
\emph{$\mathcal G$-covariant} if $E[\phi;b]$ is itself a covariant quantity
built from $(\phi,b)$, namely if there is a representation $D_E$ with
\begin{equation}\label{eq:covariant-eq}
  E[\,g\cdot\phi;\;g\cdot b\,] \;=\; D_E(g)\,E[\phi;b]
  \qquad\text{for all } g\in\mathcal G \text{ and all } (\phi,b).
\end{equation}
\end{definition}

Equation~\eqref{eq:covariant-eq} is the precise form of what older texts call
\emph{form invariance}: the same operator $E$, evaluated on the transformed
fields \emph{and} transformed background, is the transformed original. Note that
$g$ acts on $b$ as well; this is essential below.

\begin{definition}[Symmetry of the equations]\label{def:symmetry}
A transformation $g$ is a \emph{symmetry} of the system \eqref{eq:EOM} at
fixed background $b_0$ if it maps solutions to solutions:
$E[\phi;b_0]=0\ \Rightarrow\ E[g\cdot\phi;\,b_0]=0$.
\end{definition}

\begin{proposition}[Covariance $+$ invariance of the background $=$ symmetry]
\label{prop:main}
Let the equations be $\mathcal G$-covariant in the sense of
\eqref{eq:covariant-eq}, and let $b_0$ be a given background. Then the
stabilizer
\begin{equation}\label{eq:stabilizer}
  G \;=\; \Stab_{\mathcal G}(b_0)\;=\;\{\,g\in\mathcal G\;:\; g\cdot b_0=b_0\,\}
\end{equation}
is a symmetry group of $E[\phi;b_0]=0$.
\end{proposition}
\begin{proof}
For $g\in G$ and a solution $\phi$,
$E[g\cdot\phi;b_0]=E[g\cdot\phi;g\cdot b_0]=D_E(g)\,E[\phi;b_0]=0.$
\end{proof}

This is the formal statement of \eqref{eq:slogan}. Three different things
can be varied independently: the group $\mathcal G$ (how generous the
formulation is), the background $b_0$ (what is declared non-dynamical), and
the resulting physical symmetry group $G$.

\begin{remark}[Kretschmann's observation]\label{rem:kretschmann}
Covariance alone constrains nothing. Any equation can be made covariant under
\emph{any} group by promoting whatever breaks the covariance to a background
field and letting it transform along. Maxwell's equations written in
curvilinear coordinates, $\nabla_\mu F^{\mu\nu}=\tfrac{4\pi}{c}j^\nu$ with
the Christoffel symbols of the flat metric, are covariant under all
diffeomorphisms of $M$ provided $\eta_{\mu\nu}(x)$ is transformed as a
tensor field; yet the stabilizer of $\eta$ is still just the Poincar\'e
group. The physical content sits entirely in the stabilizer \eqref{eq:stabilizer},
i.e.\ in what the background is and whether it is invariant. This was pointed
out by Kretschmann in 1917 as an objection to Einstein's claim that general
covariance had physical significance \cite{kretschmann,norton}; the resolution
is discussed in Sections~\ref{sec:gr} and~\ref{sec:kretschmann}.
\end{remark}

\begin{remark}[Larger covariance, smaller symmetry]\label{rem:smaller}
Enlarging $\mathcal G$ does not enlarge the symmetry group $G$. It only
enlarges the list of structures that must be declared as backgrounds, and $G$
is whatever preserves them. This is the ``absolute objects'' criterion of
Anderson \cite{anderson}: the symmetry group of a theory is the subgroup of
the covariance group that preserves its absolute (non-dynamical) objects.
\end{remark}

\begin{definition}[Variational symmetry]\label{def:variational}
Let the equations derive from an action, $E=\delta S/\delta\phi$. A
transformation $g$ is a \emph{symmetry of the action} (a variational
symmetry) if $S[g\cdot\phi;b_0]=S[\phi;b_0]$ up to boundary terms, i.e.\ if
the Lagrangian density changes by a total divergence.
\end{definition}

\begin{proposition}\label{prop:variational}
A variational symmetry is a symmetry of the equations. The converse is false.
\end{proposition}
\begin{proof}[Sketch]
Differentiate $S[g\cdot\phi]=S[\phi]$ in an arbitrary direction $\delta\phi$:
$\langle E[g\cdot\phi],\,Dg\,\delta\phi\rangle=\langle E[\phi],\delta\phi\rangle$,
where $Dg$ is the (invertible) linearization of $g$. Hence
$E[g\cdot\phi]=(Dg^{\dagger})^{-1}E[\phi]$, so $E[\phi]=0$ implies
$E[g\cdot\phi]=0$. For the converse, see Section~\ref{sec:summary} (rescaling
a linear equation; electromagnetic duality).
\end{proof}

Covariance of the \emph{action} means that $S$ is a scalar (the Lagrangian
is a scalar density) under $\mathcal G$ acting on $(\phi,b)$; together with
invariance of $b_0$ it yields a variational symmetry, which is the entry
point for Noether's theorem (Section~\ref{sec:noether}). Figure~\ref{fig:logic}
sketches the logical architecture.

\begin{figure}[ht]
\centering
\begin{tikzpicture}[
  box/.style={draw, rounded corners=2pt, align=center, inner sep=5pt,
              font=\small, text width=4.1cm},
  arr/.style={-{Stealth[length=2.4mm]}, thick}]
  \node[box] (cov) at (-3.6,0)    {Covariant equations\\ $E[\phi;b]=0$ under $\mathcal G$};
  \node[box] (inv) at ( 3.6,0)    {Invariant background\\ $g\cdot b_0=b_0$};
  \node[box] (sym) at (0,-2.4)    {Symmetry group\\ $G=\Stab_{\mathcal G}(b_0)$};
  \node[box,text width=5cm] (n1)  at (-3.6,-4.9) {Global symmetry of the action\\ $\Rightarrow$ conserved currents\\
                                   (Noether I)};
  \node[box,text width=5cm] (n2)  at ( 3.6,-4.9) {Arbitrary functions of $x$ (no background)\\
                                   $\Rightarrow$ gauge redundancy, identities\\ (Noether II)};
  \draw[arr] (cov) -- (sym);
  \draw[arr] (inv) -- (sym);
  \draw[arr] (sym) -- (n1);
  \draw[arr] (sym) -- (n2);
\end{tikzpicture}
\caption{Logical structure of Proposition~\ref{prop:main} and its
consequences.}
\label{fig:logic}
\end{figure}

% ---------------------------------------------------------------------
\section{The infinitesimal form}
\label{sec:infinitesimal}

Let $\xi=\xi^\mu\partial_\mu$ generate a one-parameter subgroup $g_\eps$.
For a tensor field $T$ the active action to first order is
\begin{equation}\label{eq:lie}
  g_\eps\cdot T = T-\eps\,\Lie_\xi T+O(\eps^2),
\end{equation}
where $\Lie_\xi$ is the Lie derivative; for a scalar,
$\delta\phi=-\eps\,\xi^\mu\partial_\mu\phi$, reproducing
\eqref{eq:active}. The three notions become:
\begin{description}[leftmargin=1.2em,itemsep=3pt]
\item[Invariance of an object] $\Lie_\xi T=0$. For the background metric,
  $\Lie_\xi\eta_{\mu\nu}=\partial_\mu\xi_\nu+\partial_\nu\xi_\mu=0$ is the
  Killing equation. Its general solution
  $\xi^\mu=a^\mu+\omega^\mu{}_{\nu}x^\nu$, $\omega_{\mu\nu}=-\omega_{\nu\mu}$,
  has $4+6=10$ parameters: the Poincar\'e algebra.
\item[Covariance of an equation] $E[\phi;b]$ is a tensor field built
  \emph{naturally} from $\phi$ and $b$, so that
  \begin{equation}\label{eq:naturality}
    \Lie_\xi E[\phi;b]=E'[\phi;b]\cdot\Lie_\xi\phi+\frac{\partial E}{\partial b}\cdot\Lie_\xi b,
  \end{equation}
  with $E'$ the linearization of $E$ in $\phi$.
\item[Symmetry] If $\Lie_\xi b_0=0$ and $\phi$ is a solution, then
  $E[\phi;b_0]=0$ implies $\Lie_\xi E=0$, and \eqref{eq:naturality} reduces to
  \begin{equation}\label{eq:zero-mode}
    E'[\phi;b_0]\cdot\Lie_\xi\phi=0.
  \end{equation}
  Thus $\Lie_\xi\phi$ solves the \emph{linearized} equations: a symmetry
  generates a one-parameter family of new solutions, $\phi_\eps=\exp(-\eps\Lie_\xi)\phi$.
\end{description}

% ---------------------------------------------------------------------
\section{Symmetry of a law versus symmetry of a solution}
\label{sec:law-solution}

Definition~\ref{def:symmetry} concerns the \emph{law}. A separate and equally
important notion concerns an individual \emph{solution}.

\begin{definition}\label{def:isotropy}
The \emph{isotropy group} (or stabilizer) of a solution $\phi_0$ is
$G_{\phi_0}=\{g\in G: g\cdot\phi_0=\phi_0\}$, the subgroup of the symmetry group
under which \emph{that solution} is invariant. The set $G\cdot\phi_0
\cong G/G_{\phi_0}$ is the orbit of solutions generated from $\phi_0$.
\end{definition}

The distinction is the source of many apparent paradoxes:
\begin{itemize}[leftmargin=1.6em,itemsep=2pt]
\item The static Coulomb field of a point charge is invariant under rotations
  about the charge, $G_{\phi_0}=SO(3)\subset G$, although the law is invariant
  under the full Poincar\'e group. Translations generate the orbit of all
  Coulomb fields, one for each position of the charge.
\item A circular Kepler orbit in a central potential is invariant under
  neither time translations nor rotations about the orbital axis separately,
  but under the \emph{combined} one-parameter subgroup generated by
  $\partial_t+\omega\,\partial_\varphi$: a helical symmetry. Elliptic orbits
  have no continuous isotropy at all, although the law has $SO(3)\times
  \mathbb R_t$ (and more, Section~\ref{sec:kepler}).
\item In general relativity (Section~\ref{sec:gr}) the \emph{law} is invariant
  under all diffeomorphisms, but the \emph{solutions} typically have few or no
  Killing vectors; conserved quantities exist only for the symmetries of a
  given solution.
\item In spontaneous symmetry breaking the vacuum solution has a smaller
  isotropy group than the law (Section~\ref{sec:kepler}).
\end{itemize}
Symmetry of the law is a property of the theory; symmetry of the solution is a
property of a state (initial and boundary conditions included). Conservation
laws follow from the former, \emph{selection rules} and simplifying ans\"atze
from the latter.

% ---------------------------------------------------------------------
\section{Noether's theorems}
\label{sec:noether}

\subsection{First theorem: global symmetries}

Consider a Lagrangian density $\mathcal L(\phi_a,\partial_\mu\phi_a)$ with
the background $b_0$ suppressed, and an infinitesimal transformation
$\delta\phi_a=\eps X_a$ under which
\begin{equation}\label{eq:quasi}
  \delta\mathcal L=\eps\,\partial_\mu K^\mu
\end{equation}
(a variational symmetry in the sense of Definition~\ref{def:variational};
$K^\mu=0$ if $\mathcal L$ is strictly invariant). Writing $\delta\mathcal L$
in the standard way,
\begin{equation}
  \delta\mathcal L
  =\Big[\pd{\mathcal L}{\phi_a}-\partial_\mu\pd{\mathcal L}{(\partial_\mu\phi_a)}\Big]\delta\phi_a
   +\partial_\mu\Big(\pd{\mathcal L}{(\partial_\mu\phi_a)}\,\delta\phi_a\Big),
\end{equation}
and comparing with \eqref{eq:quasi} gives, with $\mathrm{EL}_a$ the
Euler--Lagrange expressions,
\begin{equation}\label{eq:noether1}
  \partial_\mu J^\mu=-\,\mathrm{EL}_a X_a\;\weakeq\;0,
  \qquad
  J^\mu=\pd{\mathcal L}{(\partial_\mu\phi_a)}X_a-K^\mu .
\end{equation}
Here $\weakeq$ means equality on solutions. The charge
$Q=\int J^0\,\dd^3x$ is conserved if the fields fall off at spatial infinity.

\begin{example}[Translations]\label{ex:translations}
A translation $x^\mu\to x^\mu+a^\mu$ acts as $\delta\phi_a=-a^\nu\partial_\nu\phi_a$.
If $\mathcal L$ has no explicit dependence on $x$ then
$\delta\mathcal L=-a^\nu\partial_\nu\mathcal L=\partial_\mu(-a^\mu\mathcal L)$,
so $K^\mu=-a^\mu\mathcal L$ and \eqref{eq:noether1} gives the canonical
energy--momentum tensor,
\begin{equation}
  \partial_\mu T^\mu{}_\nu\weakeq0,\qquad
  T^\mu{}_\nu=\pd{\mathcal L}{(\partial_\mu\phi_a)}\,\partial_\nu\phi_a-\delta^\mu_\nu\,\mathcal L,
\end{equation}
with $T^0{}_0$ the energy density. If $\mathcal L$ depends explicitly on $x$,
which means precisely that a non-invariant background is present, then
$\partial_\mu T^\mu{}_\nu\weakeq-\,(\partial\mathcal L/\partial x^\nu)_{\rm expl}$
and the corresponding component is not conserved. Proposition~\ref{prop:main}
shows up here in the guise of Noether's theorem.
\end{example}

Table~\ref{tab:noether} lists the standard pairs of symmetry and conserved
quantity. Note that conservation requires a \emph{variational} symmetry; a
symmetry of the equations that is not variational produces no Noether charge
in this way.

\begin{table}[ht]
\centering
\small
\renewcommand{\arraystretch}{1.2}
\begin{tabular}{@{}lcl@{}}
\toprule
Symmetry & Parameters & Conserved quantity\\
\midrule
time translation & 1 & energy\\
space translations & 3 & momentum\\
rotations & 3 & angular momentum\\
Galilean boosts (quasi-invariance) & 3 & centre-of-mass motion, $m\vx-\vp\,t$\\
Lorentz boosts & 3 & boost charge $J^{0i}$\\
global phase $\phi\to e^{\ii\alpha}\phi$ & 1 & electric charge\\
dilatations (conformal theories) & 1 & dilatation charge\\
\bottomrule
\end{tabular}
\caption{Some symmetry--conservation pairs (Noether's first theorem).}
\label{tab:noether}
\end{table}

\subsection{Second theorem: local symmetries}

Suppose the action is invariant under transformations that depend on
\emph{arbitrary functions} $\eps^\alpha(x)$,
$\delta\phi_a=R_{a\alpha}\eps^\alpha$ with $R$ a differential operator. Then
$0=\delta S=\int \mathrm{EL}_a R_{a\alpha}\eps^\alpha\,\dd^4x$ for all
$\eps^\alpha$, and after integrating by parts
\begin{equation}\label{eq:noether2}
  R^{\dagger}_{\alpha a}\,\mathrm{EL}_a\;\equiv\;0\qquad\text{(identically, off shell).}
\end{equation}
These \emph{Noether identities} are relations among the equations of motion
themselves, so the equations are not independent and part of the solution is
undetermined: gauge freedom.

\begin{example}[Electrodynamics]\label{ex:noether2-em}
Take
\begin{equation}\label{eq:maxwell-action}
  S=\frac1c\int \dd^4x\Big[-\frac{1}{16\pi}F_{\mu\nu}F^{\mu\nu}-\frac1c\,j^\mu A_\mu\Big],
  \quad F_{\mu\nu}=\partial_\mu A_\nu-\partial_\nu A_\mu,
\end{equation}
with $\dd^4x=c\,\dd t\,\dd^3x$. Varying $A_\nu$ gives
$\partial_\mu F^{\mu\nu}=\tfrac{4\pi}{c}j^\nu$. The transformation
$\delta A_\mu=\partial_\mu\Lambda(x)$ leaves the field part invariant; the
Noether identity is $\partial_\nu\,\mathrm{EL}^\nu\equiv0$, i.e.\
$\partial_\nu\partial_\mu F^{\mu\nu}\equiv0$ on the field side. Consistency of
the field equations then demands $\partial_\nu j^\nu=0$: charge conservation
appears as a \emph{consistency condition}, not as the charge of a
global symmetry of the field.
\end{example}

\begin{remark}[Global symmetry versus gauge redundancy]\label{rem:gauge}
Global symmetries (finitely many constant parameters) map a configuration to a
\emph{physically different} one that is also a solution; they have conserved
charges that can be measured. Local transformations map a configuration to
a \emph{physically identical} one, written differently; they express
redundancy of the description, not a relation between distinct physical
situations. ``Gauge symmetry'' is thus a symmetry of the formulation in the
sense of covariance, and the word symmetry is used for it only by courtesy.
\end{remark}

% ---- SUBSECTION BEGINS ----
\subsection{Finite and infinite numbers of degrees of freedom}
\label{sec:dof}

The chapter moves freely between mechanical systems (the Galilei group, the
Kepler problem, the relativistic particle) and field theories
(electrodynamics, general relativity). This is legitimate because the two are
not different theories: a mechanical system is a field theory on a
one-dimensional base manifold, and a field theory is a mechanical system whose
coordinates carry a continuous label. This subsection makes both statements
explicit and notes where the reduction fails to be a pure special case.

\paragraph{Mechanics as a field theory on $M=\mathbb R_t$.}
Take the base manifold to be the time axis, with the single coordinate
$x^0\equiv t$, and identify
\begin{equation}\label{eq:dictionary-id}
  \varphi_a(x^\mu)\;\leftrightarrow\;q^i(t),\qquad
  \mathcal L(\varphi_a,\partial_\mu\varphi_a;x)\;\leftrightarrow\;L(q^i,\dot q^i,t),\qquad
  \partial_\mu\;\leftrightarrow\;\frac{\dd}{\dd t}.
\end{equation}
The formulas of Sections~\ref{sec:definitions} and~\ref{sec:noether} then
specialize as in Table~\ref{tab:dictionary-dof}.

\begin{table}[ht]
\centering
\small
\renewcommand{\arraystretch}{1.35}
\setlength{\tabcolsep}{4pt}
\begin{tabular}{@{}>{\raggedright\arraybackslash}p{7.3cm}>{\raggedright\arraybackslash}p{7.3cm}@{}}
\toprule
Field theory on $M^d$ & Mechanics ($d=1$)\\
\midrule
Fields $\varphi_a(x^\mu)$ & Coordinates $q^i(t)$\\
Euler--Lagrange: $\pd{\mathcal L}{\varphi_a}-\partial_\mu\pd{\mathcal L}{(\partial_\mu\varphi_a)}=0$
  & $\pd{L}{q^i}-\dfrac{\dd}{\dd t}\pd{L}{\dot q^i}=0$\\
Noether current: $J^\mu=\pd{\mathcal L}{(\partial_\mu\varphi_a)}X_a-K^\mu$,\ \ $\partial_\mu J^\mu\weakeq0$
  & Noether charge: $Q=p_iX^i-K$,\ \ $\dot Q\weakeq0$\\
Energy--momentum tensor: $T^\mu{}_\nu=\pd{\mathcal L}{(\partial_\mu\varphi_a)}\partial_\nu\varphi_a-\delta^\mu_\nu\mathcal L$
  & $T^0{}_0=p_i\dot q^i-L=H$\\
Explicit $x$-dependence: $\partial_\mu T^\mu{}_\nu\weakeq-\big(\pd{\mathcal L}{x^\nu}\big)_{\!\rm expl}$
  & $\dot H=-\pd{L}{t}$\\
\bottomrule
\end{tabular}
\caption{Mechanics as the case $d=1$ of field theory.}
\label{tab:dictionary-dof}
\end{table}

In $d=1$ the divergence $\partial_\mu J^\mu$ is a single derivative, so the
current \emph{is} the charge and the conservation law is constancy in time.
The derivation of Noether's first theorem is the same line in both cases. For
$\delta q^i=\eps X^i$ with $\delta L=\eps\dot K$,
\begin{equation}\label{eq:noether-mech}
  \delta L=\Big[\pd{L}{q^i}-\frac{\dd}{\dd t}\pd{L}{\dot q^i}\Big]\eps X^i
  +\eps\,\frac{\dd}{\dd t}\big(p_iX^i\big)=\eps\dot K
  \;\Longrightarrow\;
  \frac{\dd}{\dd t}\big(p_iX^i-K\big)\weakeq0,\qquad p_i=\pd{L}{\dot q^i},
\end{equation}
which is \eqref{eq:noether1} with one base coordinate. The Galilean boost of
Section~\ref{sec:galilei}, $X^i=u^it$ and $K=m\,\bm u\!\cdot\!\vx$, is an
application of exactly this formula.

\paragraph{The covariance machinery in $d=1$.}
Proposition~\ref{prop:main} applies with $\mathcal G=\Diff(\mathbb R)$, the
reparametrizations of time, and with the preferred chart, that is the one-form
$\dd t$, as the background $b_0$. A reparametrization $t\to f(t)$ pulls $\dd t$
back to $f'(t)\,\dd t$, so
\begin{equation}
  \Stab_{\Diff(\mathbb R)}(\dd t)=\{\,t\to t+c\,\}.
\end{equation}
The symmetry is time translation, and Noether's theorem gives conservation of
$H$. A prescribed time-dependent external force is a second background; it
shrinks the stabilizer to the identity and produces $\dot H=-\partial L/\partial t$,
the last row of Table~\ref{tab:dictionary-dof}. Noether's \emph{second}
theorem applies to the parametrized action of Example~\ref{ex:param} (or of
Section~\ref{sec:sr}), where $t$ is promoted to a dynamical variable: the
identity \eqref{eq:param-identity} is \eqref{eq:noether2} for
$\Diff(\mathbb R)$. The soldering form of Section~\ref{sec:soldering} reduces
to the single function $e_\lambda=\dd X^0/\dd\lambda=\dot t$, with
$D=\dot t^{-1}\dd/\dd\lambda$; the parametrized Newtonian action is
$\int\dot t\,\tfrac m2(D\vx)^2\,\dd\lambda$, which equals $\int\tfrac m2\dot\vx^2/\dot t\,\dd\lambda$.

\paragraph{Base symmetries and target symmetries.}
A mechanical system is a map $q\colon M^1\to N$ from the time axis into a
\emph{target} space $N$ (a sigma model). Transformations then split into those
acting on the base, which change the arguments of the fields, and those acting
on the target, which change their values. Table~\ref{tab:base-target} shows
that the distinction covers every example of the chapter.

\begin{table}[ht]
\centering
\footnotesize
\renewcommand{\arraystretch}{1.3}
\setlength{\tabcolsep}{3pt}
\begin{tabular}{@{}>{\raggedright\arraybackslash}p{3.0cm}%
                  >{\raggedright\arraybackslash}p{2.1cm}%
                  >{\raggedright\arraybackslash}p{2.4cm}%
                  >{\raggedright\arraybackslash}p{3.0cm}%
                  >{\raggedright\arraybackslash}p{3.6cm}@{}}
\toprule
System & Base $M$ & Target / fibre & Base symmetries & Target symmetries\\
\midrule
Newtonian particle & $\mathbb R_t$ & $E^3$ & time translation
  & $E(3)$; boosts are target maps with $t$-dependent parameter\\
Relativistic particle & $\mathbb R_\lambda$ & Minkowski space
  & $\Diff(\mathbb R)$ (gauge) & Poincar\'e\\
Real scalar field & Minkowski space & $\mathbb R$ & Poincar\'e
  & $\varphi\to-\varphi$ if $V$ is even\\
$O(N)$ scalar model & Minkowski space & $\mathbb R^N$ & Poincar\'e & $O(N)$\\
Gravity (tetrad form) & $M^4$ & fields $e^a{}_\mu$ & $\Diff(M)$ (gauge)
  & local Lorentz (gauge)\\
\bottomrule
\end{tabular}
\caption{Base and target symmetries.}
\label{tab:base-target}
\end{table}

\paragraph{Field theory as a mechanical system with a continuum of degrees of
freedom.}
In the opposite direction, a $3{+}1$ split of a field theory turns it into a
mechanical system: set $q_{\vx}(t)=\varphi(\vx,t)$, so that the discrete index
$i$ becomes the continuous label $\vx$. Sums become integrals,
\begin{equation}
  L=\int\dd^3x\,\mathcal L,\qquad
  \pi(\vx)=\frac{\delta L}{\delta\dot\varphi(\vx)},\qquad
  Q=\sum_ip_iX^i-K\ \longrightarrow\ \int\dd^3x\,\big(\pi X-K^0\big)=\int\dd^3x\,J^0 .
\end{equation}
A chain of masses shows how the two descriptions relate. On a ring of $N$
masses,
\begin{equation}
  L=\sum_{n}\Big[\tfrac12m\,\dot q_n^2-\tfrac12k\,(q_{n+1}-q_n)^2\Big].
\end{equation}
Set $q_n(t)=\varphi(na,t)$, $m=\rho a$, $k=\tau/a$, and keep the length $Na$
fixed. As $a\to0$,
\begin{equation}
  L\ \to\ \int\dd x\,\Big[\tfrac12\rho\,\dot\varphi^2-\tfrac12\tau\,\varphi'^2\Big],
  \qquad\rho\ddot\varphi=\tau\varphi'' .
\end{equation}
The same system has three different symmetry statements.
\begin{enumerate}[leftmargin=1.8em,itemsep=3pt,label=\textbf{\arabic*.}]
\item \emph{Shift} $q_n\to q_n+\eps$ for all $n$, a target symmetry with $X=1$,
  $K=0$. The charge $\sum_np_n=\sum_nm\dot q_n$ becomes $\int\rho\dot\varphi\,\dd x$,
  with the local current $J^t=\rho\dot\varphi$, $J^x=-\tau\varphi'$, and
  $\partial_tJ^t+\partial_xJ^x=\rho\ddot\varphi-\tau\varphi''=0$.
\item \emph{Lattice translation} $n\to n+1$, a discrete symmetry (the cyclic group
  $\mathbb Z_N$ of the ring). It has no Noether charge in the sense of
  Section~\ref{sec:noether}.
\item \emph{Continuum translation} $x\to x+b$, a base symmetry that exists only
  in the limit. By \eqref{eq:noether1} with $\delta\varphi=-b\,\partial_x\varphi$
  its charge is the field momentum $P=-\int\rho\,\dot\varphi\,\varphi'\,\dd x$.
\end{enumerate}
So the relation between the finite and the infinite system is not only
``special case.'' The finite chain has \emph{fewer} symmetries than its
continuum limit, because discretization generically breaks continuous base
symmetries (and, for a relativistic theory on a lattice, Lorentz invariance).

\paragraph{Where the reduction is not a pure special case.}
\begin{itemize}[leftmargin=1.6em,itemsep=2pt]
\item \emph{Boundary conditions.} For $d=1$ the only boundaries are the initial
  and final times, and $Q$ needs no fall-off condition. For $d>1$ the charge
  $Q=\int J^0\,\dd^3x$ is conserved only if the flux through spatial infinity
  vanishes, which requires conditions on the fields.
\item \emph{Constraints.} A gauge parameter $\eps(\vx,t)$ yields one constraint
  per spatial point (Gauss's law, $\nabla\!\cdot\vE=4\pi\rho$, in
  Example~\ref{ex:noether2-em}), whereas reparametrization in $d=1$ yields a
  single constraint.
\item \emph{Soldering form.} It has $d$ components ($d=4$ for spacetime), not one.
\item \emph{Manifest covariance.} A $3{+}1$ split gives up manifest covariance,
  and boosts then act nonlinearly, exactly as for the gauge-fixed particle
  $\lambda=t$ of Section~\ref{sec:sr}.
\item \emph{Symmetry breaking in the quantum theory.} With finitely many
  degrees of freedom the Stone--von Neumann theorem makes all irreducible
  representations of the canonical commutation relations unitarily
  equivalent, so a symmetry is implemented unitarily and a degenerate classical
  minimum is resolved by tunneling (the ground state of a symmetric double well
  is symmetric). With infinitely many degrees of freedom inequivalent
  representations exist, and spontaneous symmetry breaking of the quantum
  vacuum (Section~\ref{sec:kepler}) becomes possible. Classically, the
  distinction between symmetry of the law and symmetry of the solution
  (Section~\ref{sec:law-solution}) is already present in both cases.
\end{itemize}
In both directions Proposition~\ref{prop:main} is unchanged: the symmetry group
is the stabilizer of the background in the covariance group. Only the
dimension of the base and the supply of target symmetries differ.
% ---- SUBSECTION ENDS ----

% ---------------------------------------------------------------------
\section{Example I: Newtonian mechanics and the Galilei group}
\label{sec:galilei}

\subsection{Covariance and invariance of Newton's law}

Take $N$ particles with $m_a\ddot{\vx}_a=\sum_{b\ne a}\bm F_{ab}$, where
$\bm F_{ab}=-\nabla_{a}V(|\vx_a-\vx_b|)$ is a central force. The Galilei group
(10 parameters) acts as
\begin{equation}\label{eq:galilei}
  t'=t+s,\qquad \vx'=R\,\vx+\vv\,t+\bm a,\qquad R\in SO(3).
\end{equation}
For \eqref{eq:galilei},
\begin{equation}
  \ddot{\vx}'_a=R\,\ddot{\vx}_a,\qquad
  \vx'_a-\vx'_b=R\,(\vx_a-\vx_b).
\end{equation}
The acceleration is thus \emph{invariant} under boosts and translations and
\emph{covariant} under rotations, and boosts and translations cancel in
differences of positions. Hence $\bm F'_{ab}=R\bm F_{ab}$ and the equation of
motion is covariant. The velocity $\dot\vx'=R\dot\vx+\vv$, by contrast, is
\emph{not} covariant under boosts (it transforms inhomogeneously), so it can
enter a boost-symmetric force law only through differences of velocities.

\subsection{Breaking by a background: the medium}

Add friction against a medium with prescribed velocity field $\bm w(\vx,t)$:
\begin{equation}\label{eq:friction}
  m\ddot\vx=-\gamma\,(\dot\vx-\bm w).
\end{equation}
Under a boost, $\bm w\to\bm w'=\bm w+\vv$, so \eqref{eq:friction} is
\emph{covariant} under the Galilei group provided $\bm w$ is transformed
along with $\vx$. But the boost is a symmetry only if the background is
unchanged. For a medium at rest ($\bm w=0$) the stabilizer consists of time
translations, space translations and rotations: boosts are excluded. The medium
plays the role of Anderson's absolute object, and its rest frame is a
preferred frame. Including the medium's atoms as dynamical degrees of
freedom restores the boost symmetry of the \emph{closed} system, exactly as in
Example~\ref{ex:coulomb}.

\subsection{Boosts: invariance up to a total derivative}

The Lagrangian $L=\sum\tfrac12m_a\dot\vx_a^2-V$ is not invariant under boosts
but is \emph{quasi-invariant}: for the infinitesimal boost
$\delta\vx_a=\eps\,\bm u\,t$,
\begin{equation}
  \delta L=\eps\sum_a m_a\bm u\cdot\dot\vx_a=\eps\,\frac{\dd}{\dd t}\Big(\bm u\cdot\sum_a m_a\vx_a\Big),
\end{equation}
which has the form \eqref{eq:quasi} with $K=\bm u\cdot\sum m_a\vx_a$. The
Noether charge $Q=\sum_a\vp_a\!\cdot\delta\vx_a/\eps-K=-\bm u\cdot\big(M\bm X-\bm P\,t\big)$
shows that
\begin{equation}
  \bm G=M\bm X_{\rm cm}-\bm P\,t=\text{const.}
\end{equation}
which says that the centre of mass moves uniformly. The quasi-invariance has a
deeper consequence: $\Poi{G_i}{P_j}=M\delta_{ij}$, so the total mass $M$
appears as a \emph{central charge} of the Galilei algebra, in contrast to
the Poincar\'e algebra where $[K_i,P_j]\propto\delta_{ij}H$. This is why mass
is a conserved parameter of Galilean systems and why its quantum version
requires a projective representation (Section~\ref{sec:quantum}).

\subsection{Non-inertial frames and the geometric formulation}

Under a time-dependent change of frame $\vx'=R(t)\vx+\bm a(t)$ Newton's law
acquires extra terms (Coriolis, centrifugal, Euler, and translational
inertial forces) and no longer has the same form. In the language of
Section~\ref{sec:definitions} the law is covariant under the Galilei group
\eqref{eq:galilei} only. It can, however, be \emph{reformulated} so that it is
covariant under arbitrary coordinate changes: this is Cartan's
Newton--Cartan geometry \cite{cartan,dbkp}. Free fall is the geodesic equation
of a connection with $\Gamma^i{}_{00}=\partial_i\Phi$, and the absolute
objects are the time one-form $t_\mu=\partial_\mu t$ and the spatial
cometric $h^{\mu\nu}$. The formulation is generally covariant, the symmetry
is the subgroup preserving $t_\mu$ and $h^{\mu\nu}$, which is larger than the
Galilei group (it includes $\vx\to\vx+\bm a(t)$ acting on $\Phi$ by a term
linear in $\vx$, proportional to $\ddot{\bm a}$, an expression of the
Newtonian equivalence principle). A textbook illustration of
Remark~\ref{rem:kretschmann}.

% ---------------------------------------------------------------------
\section{Example II: electrodynamics, special relativity and the point particle}
\label{sec:sr}

\subsection{Lorentz covariance of Maxwell's equations}

With $F^{0i}=-E^i$ and $F^{ij}=-\epsilon^{ijk}B^k$, the Maxwell equations in
Gaussian units read
\begin{equation}\label{eq:maxwell}
  \partial_\mu F^{\mu\nu}=\frac{4\pi}{c}\,j^\nu,\qquad
  \partial_\mu\tilde F^{\mu\nu}=0,\qquad
  \tilde F^{\mu\nu}=\tfrac12\epsilon^{\mu\nu\alpha\beta}F_{\alpha\beta}.
\end{equation}
Under a Lorentz transformation \eqref{eq:poincare}, $F^{\mu\nu}$ transforms as a
rank-two tensor and $j^\nu$ as a vector, so
$\mathcal E^\nu\equiv\partial_\mu F^{\mu\nu}-\tfrac{4\pi}{c}j^\nu$ is a
four-vector: $\mathcal E'^\nu=\Lambda^\nu{}_\alpha\mathcal E^\alpha$. This is
\eqref{eq:covariant-eq} with $D_E(g)=\Lambda$; the metric $\eta_{\mu\nu}$,
needed to raise and lower indices, is invariant, so Proposition~\ref{prop:main}
gives Poincar\'e symmetry.

It is worth stating exactly what is covariant and what is invariant.
The fields $\vE$ and $\vB$ are \emph{components} of a tensor and mix under
boosts; for a boost with velocity $c\vbeta$,
\begin{equation}\label{eq:boost-fields}
  \vE'_\parallel=\vE_\parallel,\quad \vB'_\parallel=\vB_\parallel,\quad
  \vE'_\perp=\gamma(\vE_\perp+\vbeta\times\vB),\quad
  \vB'_\perp=\gamma(\vB_\perp-\vbeta\times\vE).
\end{equation}
The \emph{invariants} are scalars built from $F$ and $\eta$ (and $\epsilon$):
\begin{equation}\label{eq:EM-invariants}
  F_{\mu\nu}F^{\mu\nu}=2(B^2-E^2),\qquad
  F_{\mu\nu}\tilde F^{\mu\nu}=-4\,\vE\!\cdot\!\vB .
\end{equation}
A purely electric field acquires a magnetic component in a moving frame, yet
$E^2-B^2$ and $\vE\!\cdot\!\vB$ do not change. For example, a field with
$\vE\!\cdot\!\vB=0$ and $E^2>B^2$ can be boosted to a frame in which it is
purely electric, but never to one in which it is purely magnetic. The
Lagrangian density in \eqref{eq:maxwell-action} is $\tfrac1{8\pi}(E^2-B^2)$, a
Lorentz scalar: covariance of the action in the sense of
Section~\ref{sec:definitions}. (Because $\epsilon^{\mu\nu\alpha\beta}$ is a
pseudotensor, covariance under parity requires $\tilde F$ to be treated as
one.)

The same equations are \emph{not} covariant under Galilean transformations
with any local rule for $\vE,\vB$; limits in which they are (``electric'' and
``magnetic'' Galilean electromagnetism) exist, but each captures only part of
the equations \cite{lbll}. The failure of Galilean covariance of Maxwell's
theory was the original motivation for special relativity.

\subsection{Larger symmetries and symmetry without an action}

The Poincar\'e group is not the whole story.
\begin{itemize}[leftmargin=1.6em,itemsep=2pt]
\item \emph{Conformal symmetry.} In four dimensions the sourceless Maxwell
  equations are covariant under all diffeomorphisms that preserve the
  \emph{conformal class} $[\eta]$, i.e.\ $\eta\to\Omega^2\eta$; the
  corresponding symmetry group is the 15-parameter conformal group
  (translations 4, Lorentz 6, dilatation 1, special conformal 4). The
  stabilizer of $[\eta]$ is bigger than the stabilizer of $\eta$, which is why
  the symmetry is bigger: Proposition~\ref{prop:main} with a weaker background.
\item \emph{Duality.} In the absence of sources, the replacement
  $\vE\to\vE\cos\alpha+\vB\sin\alpha$, $\vB\to-\vE\sin\alpha+\vB\cos\alpha$
  maps solutions to solutions, so it is a symmetry of the \emph{equations}.
  It is not a local variational symmetry of the action \eqref{eq:maxwell-action}
  built from $A_\mu$ (for $\alpha=\pi/2$ the Lagrangian density changes sign),
  though conserved charges can be obtained in a reformulation involving a
  second potential \cite{deser}. This illustrates the failure of the converse
  in Proposition~\ref{prop:variational}.
\end{itemize}

\subsection{The relativistic point particle: reparametrization as gauge symmetry}

For a free particle, the manifestly covariant action is
\begin{equation}\label{eq:particle}
  S=-mc\int\sqrt{\eta_{\mu\nu}\,\dot x^\mu\dot x^\nu}\;\dd\lambda ,
  \qquad \dot{}\equiv\frac{\dd}{\dd\lambda}.
\end{equation}
It is invariant under two very different kinds of transformation.

\emph{(i) Poincar\'e transformations} \eqref{eq:poincare}, a global
symmetry (the Lagrangian is a scalar and $\eta$ is invariant). Noether's first
theorem gives the conserved quantities $\pi_\mu=\partial L/\partial\dot x^\mu$
and $J^{\mu\nu}=x^\mu\pi^\nu-x^\nu\pi^\mu$.

\emph{(ii) Reparametrizations} $\lambda\to f(\lambda)$, which depend on an
arbitrary function: a local symmetry, $\delta x^\mu=\eps(\lambda)\dot x^\mu$.
Noether's second theorem \eqref{eq:noether2} gives the identity
\begin{equation}\label{eq:param-identity}
  \dot x^\mu\,\mathrm{EL}_\mu\equiv0,
\end{equation}
so only three of the four equations for $x^\mu(\lambda)$ are independent, and
one function (the parametrization) is undetermined. In the Hamiltonian
description the momenta satisfy the identity
$\eta^{\mu\nu}\pi_\mu\pi_\nu-m^2c^2=0$ (a primary constraint), the canonical
Hamiltonian vanishes because $L$ is homogeneous of degree one in the
velocities, and the dynamics is generated by the constraint times an arbitrary
multiplier.

Choosing the gauge $\lambda=t$ gives
$L=-mc^2\sqrt{1-v^2/c^2}$, which is \emph{not} manifestly covariant: the
boost acts \emph{nonlinearly} on $x^i(t)$, since a boost mixes $t$ with the
coordinates and the gauge condition must be restored afterwards. The
Poincar\'e symmetry is still there; it is merely hidden by the gauge choice.
Covariance is a property of the formulation, symmetry is not.

% ---------------------------------------------------------------------
\section{Example III: general relativity}
\label{sec:gr}

\subsection{General covariance is not a symmetry of the world}

Einstein's field equations in the convention of Landau and Lifshitz
(signature $(+,-,-,-)$),
\begin{equation}\label{eq:einstein}
  R_{\mu\nu}-\tfrac12 g_{\mu\nu}R=\frac{8\pi G}{c^4}\,T_{\mu\nu},
\end{equation}
derive from $S=-\tfrac{c^3}{16\pi G}\int R\sqrt{-g}\,\dd^4x+S_{\rm matter}$
and are covariant under the full diffeomorphism group $\Diff(M)$. In
Proposition~\ref{prop:main}, $\mathcal G=\Diff(M)$ and the metric $g_{\mu\nu}$
is \emph{dynamical}: there is no fixed background $b_0$ (in the absence of
fixed matter fields). The stabilizer is therefore the entire group,
\begin{equation}
  G=\Stab_{\Diff(M)}(\varnothing)=\Diff(M),
\end{equation}
and by Remark~\ref{rem:gauge} diffeomorphisms are a \emph{gauge redundancy}
rather than a physical symmetry. Here lies the resolution of Kretschmann's
objection: Einstein's theory is not special because its equations can be
written covariantly (every theory can be), but because they contain
\emph{no absolute objects} \cite{anderson} (Section~\ref{sec:kretschmann} gives
the details). In special relativity written in arbitrary coordinates the metric $\eta$ is an absolute object; in GR it is
replaced by the dynamical $g_{\mu\nu}$.

\subsection{The hole argument}

Let $g$ be a solution on $M$, and let $\psi$ be a diffeomorphism equal to the
identity outside a compact region $H$. Then $\psi^*g$ is also a solution,
equal to $g$ outside $H$ and different inside. If points of $M$ were taken to
have physical identity (a metric field being \emph{the} object of physical
interest), determinism would fail: the Einstein equations plus data outside $H$
do not determine the field inside. The way out is to regard $g$ and $\psi^*g$
as one and the same physical situation. Observables are then
diffeomorphism-\emph{invariant} (relational) quantities: for instance,
proper times along worldlines between their intersections. This is exactly
the gauge-redundancy interpretation of Remark~\ref{rem:gauge}
\cite{earman,norton}.

\subsection{Noether's second theorem in GR: conservation only for Killing vectors}

Under $\delta g_{\mu\nu}=-(\nabla_\mu\xi_\nu+\nabla_\nu\xi_\mu)$ with arbitrary
$\xi^\mu(x)$ the action is invariant, and \eqref{eq:noether2} becomes the
contracted Bianchi identity $\nabla_\mu G^{\mu\nu}\equiv0$. Consequently the
field equations \eqref{eq:einstein} imply $\nabla_\mu T^{\mu\nu}=0$, which
\emph{is not a conservation law} in the sense of a conserved charge: in a
curved spacetime $\nabla_\mu T^{\mu\nu}=0$ does not integrate to a conserved
quantity. A conserved current is obtained only if there exists a Killing vector
$\xi$ of the particular solution under consideration:
\begin{equation}
  J^\mu=T^{\mu\nu}\xi_\nu,\qquad
  \nabla_\mu J^\mu=(\nabla_\mu T^{\mu\nu})\xi_\nu+T^{\mu\nu}\nabla_{(\mu}\xi_{\nu)}=0.
\end{equation}
This is the distinction of Section~\ref{sec:law-solution} in its
sharpest form: the \emph{law} has the symmetry $\Diff(M)$ (gauge), but
conservation of energy and momentum rests on the symmetry of the \emph{solution}
(the isotropy group $G_{\phi_0}$ of the background spacetime, such as time
translation invariance of the Schwarzschild solution).

% ---------------------------------------------------------------------
\section{Kretschmann's observation in detail}
\label{sec:kretschmann}

Remark~\ref{rem:kretschmann} stated Kretschmann's observation in one
sentence. The confusion it addresses is among the most persistent in the
foundations of relativity, so we now examine it more closely: the historical
dispute, a constructive proof, what the proof does and does not show, and what
remains of ``general covariance'' afterwards.

\subsection{The dispute}

In 1915--16 Einstein presented general covariance, the form-invariance of the
field equations under arbitrary coordinate changes, as a physical principle,
for two reasons. First, it seemed to extend the principle of relativity from
inertial frames to arbitrary states of motion. Second, he argued that the
physical content of any theory is exhausted by \emph{coincidences}: the
intersection of two worldlines, the meeting of a pointer and a mark on a
scale. Coordinates merely label such coincidences, no labeling can be
physically privileged, and a theory that singles out a class of coordinates
would draw a distinction without empirical basis (the \emph{point-coincidence
argument}) \cite{norton}.

Kretschmann accepted the premise and reversed the conclusion
\cite{kretschmann}. If only coincidences are physical, he argued, then
\emph{any} spacetime theory, however un-relativistic, can be restated in a form
that does not depend on the labeling, at the price of mathematical effort: one
adds to the formalism whatever is needed to keep track of the labeling. The
demand for general covariance therefore tests the ingenuity of the theorist and
places no restriction on nature. Einstein conceded the logical point in 1918
\cite{einstein1918} but maintained that the demand retains a \emph{heuristic}
value: among generally covariant equations only the simplest and most natural
are plausible as fundamental laws, and the covariant versions of older theories
are comparatively baroque. Historically the dispute coincided with the work of
Hilbert, Klein and Noether on identities among the field equations of
generally covariant theories, and the second theorem of
Section~\ref{sec:noether} came out of that work \cite{ks}. Anderson and
Friedman later turned Kretschmann's observation into a precise criterion
\cite{anderson,friedman}, which we met in Remark~\ref{rem:smaller}; Norton's
survey \cite{norton} and his later discussion \cite{bc} are the best guides to
the debate.

\subsection{A constructive proof}

Let $E[\varphi]=0$ be \emph{any} system of differential equations for fields
$\varphi$, written in a fixed ``preferred'' chart $x^a$ on a region of
spacetime. It may be Galilean, Lorentzian, or covariant under no group at all.
In an arbitrary chart $y^\mu$ the preferred coordinates appear as four scalar
functions $X^a(y)$. Put
\begin{equation}\label{eq:soldering}
  e^a{}_\mu=\partial_\mu X^a,\qquad
  e_a{}^\mu\,e^b{}_\mu=\delta_a^b .
\end{equation}
The \emph{frame components} of a tensor field are its contractions with
$e_a{}^\mu$ and $e^a{}_\mu$ (for a covector, $\varphi_a=e_a{}^\mu\varphi_\mu$).
They carry only the index $a$ and are therefore \emph{scalars} under changes
of the chart $y$. The operators $D_a=e_a{}^\mu\partial_\mu$ map scalars to
scalars. Let $\tilde E[\tilde\varphi;e]$ be obtained from $E$ by replacing every
field by its frame components and every $\partial/\partial x^a$ by $D_a$.

\begin{proposition}[Kretschmann]\label{prop:kretschmann}
The system $\tilde E[\tilde\varphi;e]=0$ is $\Diff(M)$-covariant in the sense
of \eqref{eq:covariant-eq}, with the soldering form $e^a{}_\mu$ (Section~\ref{sec:soldering}) as background.
It is equivalent to the original system: $(\tilde\varphi,e)$ solves
$\tilde E=0$ if and only if the frame components of $\tilde\varphi$ solve
$E=0$.
\end{proposition}
\begin{proof}
Each frame component and each $D_a(\cdot)$ is a scalar under $\Diff(M)$ acting
on $(\tilde\varphi,e)$ by push-forward, so $\tilde E$ is built from scalars
and $\tilde E[\psi\!\cdot\!\tilde\varphi;\psi\!\cdot\! e]
=\psi\!\cdot\!\tilde E[\tilde\varphi;e]$ for every diffeomorphism $\psi$. In the
particular chart $y=X$ one has $e^a{}_\mu=\delta^a_\mu$ and $D_a=\partial/\partial x^a$,
so $\tilde E=E$. \qedhere
\end{proof}

Nothing about $E$ was used, which is precisely the content of Kretschmann's
claim: general covariance is available for any theory, at the cost of a
background field $e^a{}_\mu$ that the original formulation did not display.

\begin{example}[Wave equation]\label{ex:wave}
For $\eta^{ab}\partial_a\partial_b\varphi=0$ the construction gives
$\eta^{ab}D_aD_b\varphi=0$. Expanding $D_aD_b\varphi=e_a{}^\mu\partial_\mu(e_b{}^\nu\partial_\nu\varphi)$
and using $\eta^{ab}e_a{}^\mu e_b{}^\nu=\eta^{\mu\nu}$,
\begin{equation}
  \eta^{ab}D_aD_b\varphi
  =\eta^{\mu\nu}\big(\partial_\mu\partial_\nu\varphi-\Gamma^\lambda{}_{\mu\nu}\partial_\lambda\varphi\big)
  =\eta^{\mu\nu}\nabla_\mu\partial_\nu\varphi,\qquad
  \eta_{\mu\nu}=\eta_{ab}e^a{}_\mu e^b{}_\nu,
\end{equation}
where $\Gamma^\lambda{}_{\mu\nu}=e_a{}^\lambda\partial_\mu e^a{}_\nu$ are the
Christoffel symbols of $\eta_{\mu\nu}$ (flat, since $e^a{}_\mu$ is a gradient). This is exactly
how Christoffel symbols enter the ``curvilinear'' form of every equation of
special relativity.
\end{example}

\begin{example}[Newtonian mechanics, parametrized]\label{ex:param}
The free Newtonian particle, $S=\int\tfrac12m\,\dd\vx^2/\dd t$, becomes
generally covariant (here: reparametrization invariant) if $t$ is promoted to a
dynamical variable of the worldline,
\begin{equation}
  S=\int\frac m2\,\frac{\dot{\vx}^2}{\dot t}\,\dd\lambda,\qquad
  \dot{}=\frac{\dd}{\dd\lambda}.
\end{equation}
This is homogeneous of degree one in $(\dot t,\dot\vx)$, hence invariant under
$\lambda\to f(\lambda)$; choosing $\lambda=t$ gives back the original
action. The momenta obey $p_t+\vp^2/2m\weakeq0$, the Newtonian analogue of the
mass-shell constraint of Section~\ref{sec:sr}, with $-p_t$ playing the
role of the energy. This is Kretschmann's observation in one dimension.
\end{example}

\begin{remark}[Dependence on the soldering form]\label{rem:soldering-dep}
If $E$ is covariant under a group $G_0$ acting on the frame index $a$ (Lorentz
transformations, say), then $\tilde E$ depends on $e^a{}_\mu$ only through its
$G_0$-invariant parts: for the Poincar\'e group through the flat metric
$\eta_{\mu\nu}=\eta_{ab}e^a{}_\mu e^b{}_\nu$ alone, and for the Galilei group
through the Newton--Cartan data of Section~\ref{sec:galilei}.
\end{remark}

\subsection{The soldering form}
\label{sec:soldering}

The original formulation of $E$ contains two kinds of index that were never
distinguished, because they are numerically identified in the preferred chart.
The index $a$ labels the coordinates of the \emph{model space} $V=\mathbb R^4$,
on which the group $G_0$ (Poincar\'e, Galilei, \dots) acts and in which the
original theory lives. The index $\mu$ labels directions at a point of the
spacetime manifold $M$, whose charts $y^\mu$ are arbitrary. Covariantization
must pull the two apart: $\mu$ is acted on by $\Diff(M)$, whereas $a$ is merely
a label, acted on at most by $G_0$.

The one-form $e^a{}_\mu\,\dd y^\mu=\dd X^a$ is the link between them. It is a
one-form on $M$ with values in $V$: at each point $p$ it defines the linear map
\begin{equation}
  e_p\colon T_pM\to V,\qquad v^\mu\mapsto e^a{}_\mu\,v^\mu ,
\end{equation}
which is invertible precisely because $X^a$ is a chart. It identifies every
tangent space with the model space, point by point, and thereby \emph{solders}
(welds) the model space onto the manifold; hence the name. In the mathematical
literature the \emph{solder form} (or canonical form) is the tautological
$V$-valued one-form on the frame bundle, and its pull-backs by local sections
are the coframes, tetrads or vierbeins $e^a{}_\mu$ \cite{kn}. In Cartan's
generalization of Klein's Erlangen programme, a geometry is a manifold whose
tangent spaces are soldered to the model space of a Klein geometry $G_0/H$
\cite{sharpe}. (Some authors reserve ``soldering form'' for the canonical
form and call $e^a{}_\mu$ a tetrad; we follow the usage of physics.)

The soldering form does three jobs in the construction above.
\begin{enumerate}[leftmargin=1.8em,itemsep=3pt,label=\textbf{\arabic*.}]
\item \emph{It converts indices.} Frame components $\varphi_a=e_a{}^\mu\varphi_\mu$
  and the derivatives $D_a=e_a{}^\mu\partial_\mu$ are scalars under
  $\Diff(M)$ because $e^a{}_\mu$ is the only place where an index of each kind
  meets.
\item \emph{It records the preferred chart and nothing else.} Since
  $e^a{}_\mu=\partial_\mu X^a$, it is closed, $\dd e^a=0$. A translation
  $X^a\to X^a+c^a$ leaves $e^a{}_\mu$ unchanged, and a Lorentz transformation
  $X^a\to\Lambda^a{}_bX^b$ rotates it, $e\to\Lambda e$. Only the
  $G_0$-invariant content matters (Remark~\ref{rem:soldering-dep}): for the
  Poincar\'e group the metric $\eta_{\mu\nu}=\eta_{ab}e^a{}_\mu e^b{}_\nu$; for
  the Galilei group the time one-form $t_\mu=e^0{}_\mu$ (with $X^0=t$) and the
  spatial cometric $h^{\mu\nu}=\delta^{ij}e_i{}^\mu e_j{}^\nu$.
\item \emph{It defines a connection.} The coefficients
  $\Gamma^\lambda{}_{\mu\nu}=e_a{}^\lambda\partial_\mu e^a{}_\nu$ define the
  derivative for which frame components are constant; they are symmetric in
  $\mu\nu$ (no torsion) and give zero curvature, both consequences of
  $\dd e^a=0$.
\end{enumerate}

\begin{example}[Polar coordinates]\label{ex:polar}
Take the Euclidean plane with $X^1=r\cos\theta$, $X^2=r\sin\theta$. Then
\begin{equation}
  e^1=\cos\theta\,\dd r-r\sin\theta\,\dd\theta,\qquad
  e^2=\sin\theta\,\dd r+r\cos\theta\,\dd\theta,
\end{equation}
so $\delta_{ab}e^a e^b=\dd r^2+r^2\dd\theta^2$, and
$\Gamma^\lambda{}_{\mu\nu}=e_a{}^\lambda\partial_\mu e^a{}_\nu$ gives
$\Gamma^r{}_{\theta\theta}=-r$ and $\Gamma^\theta{}_{r\theta}=1/r$, the familiar
Christoffel symbols of polar coordinates. They are not extra input: they are
what the soldering form of a flat chart looks like in a curvilinear one.
\end{example}

In general relativity the same object appears as the tetrad $e^a{}_\mu$, with
local Lorentz transformations acting on $a$ as a gauge symmetry. The
difference from Kretschmann's construction lies not in the object but in its
status. There $e^a{}_\mu$ is prescribed as the gradient of a given chart, so
$\dd e^a=0$ and the metric $\eta_{ab}e^a{}_\mu e^b{}_\nu$ is flat and the same
in every solution. In general relativity $e^a{}_\mu$ is varied in the action,
$\dd e^a\neq0$ in general (the torsion-free spin connection obeys
$\dd e^a+\omega^a{}_b\wedge e^b=0$), and the metric is dynamical. The
soldering form is the fixed background of every covariantized older theory and
a dynamical field in general relativity, which is the distinction drawn by
Proposition~\ref{prop:main} and Anderson's criterion.

\subsection{Where the physics went}

After the construction the covariance group is $\Diff(M)$, acting on pairs
$(\tilde\varphi,b)$ where $b$ is the $G_0$-invariant structure of
Remark~\ref{rem:soldering-dep}. By Proposition~\ref{prop:main} the symmetry
group is $\Stab_{\Diff(M)}(b)$, and this is exactly the original $G_0$, because
$G_0$ is the automorphism group of its invariant structure. This is Klein's
Erlangen idea read backwards: a geometry determines its group, and the group
is recovered as the automorphism group of an invariant structure, as for
$\mathrm{Poincar\acute e}=\mathrm{Aut}(\eta)$ and
$\text{Galilei}=\mathrm{Aut}(t_\mu,h^{\mu\nu},\nabla)$. Covariantization
therefore changes the \emph{formulation} (the covariance group becomes
$\Diff(M)$) and leaves the \emph{physics} (the symmetry group, the solution
set) untouched. Table~\ref{tab:kretschmannize} lists the standard cases. The
friction and preferred-frame rows show what happens when the original theory
contains a genuinely non-symmetric ingredient: the covariantization turns it
into an extra background field, and the stabilizer shrinks accordingly.

\begin{table}[ht]
\centering
\small
\renewcommand{\arraystretch}{1.25}
\setlength{\tabcolsep}{4pt}
\begin{tabular}{@{}>{\raggedright\arraybackslash}p{3.7cm}%
                  >{\raggedright\arraybackslash}p{5.6cm}%
                  >{\raggedright\arraybackslash}p{5.0cm}@{}}
\toprule
Original theory & Absolute structure $b_0$ introduced & Stabilizer $\Stab_{\Diff(M)}(b_0)$\\
\midrule
Special-relativistic field theory & flat metric $\eta_{\mu\nu}$ & Poincar\'e group\\
Galilean $N$-body mechanics & $t_\mu$, $h^{\mu\nu}$, flat connection $\nabla$ & Galilei group\\
Newtonian gravity (Cartan) & $t_\mu$, $h^{\mu\nu}$; connection dynamical
  & larger than Galilei (includes $\vx\to\vx+\bm a(t)$)\\
Friction against a medium at rest & as for Galilei, plus the velocity field
  $w^\mu$ of the medium & translations, rotations; no boosts\\
Electrodynamics with a preferred frame & $\eta_{\mu\nu}$ plus a unit timelike
  vector $u^\mu$ & translations, rotations fixing $u^\mu$; no boosts\\
General relativity & none (the metric is dynamical) & $\Diff(M)$, all gauge\\
\bottomrule
\end{tabular}
\caption{Kretschmann's construction applied to several theories.}
\label{tab:kretschmannize}
\end{table}

\subsection{Two attempted escapes, and why they do not help}

\paragraph{Make the coordinates dynamical.}
One may try to remove the background by promoting the preferred chart to
dynamical fields. For a Poincar\'e-invariant action $S[\varphi;\eta]$ define the
\emph{parametrized} theory
\begin{equation}
  S_P[\varphi,X]=S\big[\varphi;\ \eta_{\mu\nu}=\eta_{ab}\,\partial_\mu X^a\partial_\nu X^b\big],
\end{equation}
now varied with respect to $X^a$ as well \cite{kuchar}. The action is
$\Diff(M)$-invariant, no chart is singled out, and by Noether's second theorem
\eqref{eq:noether2} the equations for $X^a$ are consequences of those for
$\varphi$ (they express the conservation of the energy--momentum tensor). The
$X^a$ are pure gauge: in the gauge $X^a=y^a$ one recovers special relativity.
Nothing physical has changed. By Anderson's criterion this is not a
counterexample, since the flat metric built from $X^a$ is the same in every
solution up to diffeomorphism and is therefore still an absolute object, but
it shows that the criterion must be applied to \emph{composite} objects built
from the dynamical fields as well, and that the mere presence of
diffeomorphism-gauge freedom does not distinguish general relativity from
special relativity \cite{giulini}.

\paragraph{Count the absolute objects.}
Anderson's definition calls an object \emph{absolute} if it appears in the
equations and is the same, up to diffeomorphism, in every solution, and
\emph{dynamical} otherwise. In these terms the distinguishing feature of general
relativity is that the metric is dynamical in the strong sense: different
solutions are not isometric (they differ in their Weyl curvature, and there
are gravitational waves), while every covariantized older theory carries an
inertial structure that is fixed once and for all.

\subsection{What survives of general covariance}

\begin{enumerate}[leftmargin=1.8em,itemsep=3pt,label=\textbf{\arabic*.}]
\item \textbf{As a formal requirement it is vacuous}
  (Proposition~\ref{prop:kretschmann}); it constrains the notation, not the
  theory.
\item \textbf{As gauge invariance it has consequences.} If \emph{all}
  geometric structures are dynamical, $\Diff(M)$ is a gauge group and
  Noether's second theorem gives the Bianchi identities. The consequence is
  that $\nabla_\mu T^{\mu\nu}=0$ follows from the gravitational field equations
  alone, for any matter content, so the matter equations must be compatible
  with it. In a theory with a fixed $\eta$ the same conservation law is
  instead a consequence of the matter equations. This is the part of Einstein's
  intuition that the G\"ottingen group made precise \cite{ks}.
\item \textbf{As a heuristic and a uniqueness statement it has content}, as
  Einstein argued in reply. In four dimensions the only symmetric,
  divergence-free tensor built from the metric and its first two derivatives,
  linear in the second derivatives, is $G_{\mu\nu}+\Lambda g_{\mu\nu}$
  \cite{lovelock}: combined with general covariance and a dynamical metric this
  fixes the gravitational field equations almost uniquely.
\item \textbf{The testable content lies in the stabilizer, not in
  covariance.} Universality of free fall, tested to parts in $10^{13}$ and
  beyond, and the null results of Hughes--Drever-type experiments say that
  matter couples to a single metric and that no further background structure
  (a preferred frame $u^\mu$, say) enters the equations. In the language of
  Proposition~\ref{prop:main}: they test \emph{what is held fixed}, which is the
  second summand of \eqref{eq:slogan}. The first summand, covariance, cannot be
  tested at all.
\end{enumerate}

\begin{remark}[The Kretschmann scalar]\label{rem:kretschmann-scalar}
The curvature invariant
\begin{equation}
  K=R_{\mu\nu\rho\sigma}R^{\mu\nu\rho\sigma}
\end{equation}
also carries Kretschmann's name. It illustrates the difference between covariant
\emph{components} and invariant \emph{scalars}. For the Schwarzschild solution
\begin{equation}
  K=\frac{48\,G^2M^2}{c^4\,r^6}
\end{equation}
is finite at $r=2GM/c^2$, although $g_{tt}$ vanishes and $g_{rr}$ diverges
there in Schwarzschild coordinates: the horizon is a coordinate artifact,
while $r=0$ is a genuine curvature singularity. Statements about
chart-dependent components are statements about the description; statements
about invariants are statements about the world.
\end{remark}

% ---------------------------------------------------------------------
\section{Example IV: gauge covariance}
\label{sec:gauge}

The vocabulary of Section~\ref{sec:three-words} is stretched, in a
precise way, in gauge theory. Let a charged field $\psi$ with charge $q$
transform under local $U(1)$:
\begin{equation}\label{eq:u1}
  \psi\to e^{-\ii q\Lambda(x)/\hbar c}\,\psi,\qquad
  A_\mu\to A_\mu+\partial_\mu\Lambda .
\end{equation}
The ordinary derivative $\partial_\mu\psi$ is neither invariant nor covariant:
it picks up an extra term $\propto\partial_\mu\Lambda\,\psi$. The
\emph{covariant derivative}
\begin{equation}\label{eq:covder}
  D_\mu=\partial_\mu+\frac{\ii q}{\hbar c}A_\mu
  \quad\Longrightarrow\quad
  D_\mu\psi\to e^{-\ii q\Lambda/\hbar c}\,D_\mu\psi
\end{equation}
transforms with the \emph{same phase} as $\psi$ itself: it is
gauge-\emph{covariant} in the sense of \eqref{eq:covariant-object}, with
$D(g)$ the $U(1)$ phase. Gauge-\emph{invariant} objects are built from
covariant ones by contracting the phases away:
\begin{equation}
  \psi^\dagger\psi,\qquad(D_\mu\psi)^\dagger(D^\mu\psi),\qquad
  F_{\mu\nu}=\frac{\hbar c}{\ii q}[D_\mu,D_\nu]=\partial_\mu A_\nu-\partial_\nu A_\mu .
\end{equation}
For the non-Abelian case, $\psi\to U\psi$, $D_\mu=\partial_\mu-\ii gA_\mu$
with $A_\mu=A_\mu^aT^a$, and
\begin{equation}
  A_\mu\to UA_\mu U^{-1}-\tfrac{\ii}{g}(\partial_\mu U)U^{-1},\quad
  F_{\mu\nu}=\partial_\mu A_\nu-\partial_\nu A_\mu-\ii g[A_\mu,A_\nu]\to UF_{\mu\nu}U^{-1}.
\end{equation}
$F_{\mu\nu}$ is gauge-\emph{covariant}, and $\tr F_{\mu\nu}F^{\mu\nu}$ is
gauge-\emph{invariant}. In the Abelian case $F_{\mu\nu}$ is itself invariant,
not merely covariant, because the adjoint representation of $U(1)$ is trivial.

The interplay with Noether's theorems is instructive. The \emph{global} subgroup
($\Lambda=\text{const}$) is an honest symmetry with a conserved, measurable
charge, the electric charge of Table~\ref{tab:noether}. The \emph{local}
transformations are redundancy: they yield the Noether identity
$\partial_\mu j^\mu\weakeq0$ with $j^\mu=\partial\mathcal L_{\rm matter}/\partial A_\mu$
(Example~\ref{ex:noether2-em}), and, in the Hamiltonian formalism,
a first-class constraint, Gauss's law $\nabla\!\cdot\vE=4\pi\rho$, that
generates the gauge transformations.

% ---------------------------------------------------------------------
\section{Example V: hidden symmetry and symmetry breaking}
\label{sec:kepler}

\subsection{Kepler problem: a symmetry that is not a coordinate transformation}

For $H=\vp^2/2m-k/r$ the geometric symmetry is $SO(3)\times\mathbb R_t$
(rotations and time translation), with conserved $\vL=\vx\times\vp$ and $H$.
There is, however, an additional conserved vector, the Runge--Lenz vector
\begin{equation}\label{eq:rl}
  \bm A=\vp\times\vL-mk\,\hat{\vx},\qquad \Poi{A_i}{H}=0 .
\end{equation}
The transformation it generates, $\delta\vx=\eps\,\Poi{\vx}{\vn\!\cdot\!\bm A}$,
\begin{equation}
  \delta\vx=\eps\,\big[\,2(\vn\!\cdot\!\vx)\,\vp-(\vx\!\cdot\!\vp)\,\vn-(\vn\!\cdot\!\vp)\,\vx\,\big],
\end{equation}
involves the \emph{momentum}: it is not a point transformation of
configuration space and is not related to any covariance of the
equation $m\ddot\vx=-k\vx/r^3$ under coordinate changes. It is nonetheless a
canonical symmetry of the Hamiltonian, and the Poisson algebra
\begin{equation}
  \Poi{L_i}{L_j}=\epsilon_{ijk}L_k,\quad
  \Poi{L_i}{A_j}=\epsilon_{ijk}A_k,\quad
  \Poi{A_i}{A_j}=-2mH\,\epsilon_{ijk}L_k
\end{equation}
closes on $so(4)$ for $H<0$ (after the rescaling $\bm A\to\bm A/\sqrt{-2mH}$),
$e(3)$ for $H=0$, and $so(3,1)$ for $H>0$. The extra symmetry explains why
bound orbits are closed (the perihelion does not precess) \cite{goldstein}
and, in the quantum problem, the extra degeneracy of the hydrogen levels.
It is an example of a \emph{dynamical} (hidden) symmetry:
a symmetry of the law that is not visible as a covariance of the formulation.

\subsection{Spontaneous symmetry breaking}

An $O(N)$-invariant scalar theory,
$\mathcal L=\tfrac12\partial_\mu\phi_a\partial^\mu\phi_a-\tfrac\lambda4(\phi_a\phi_a-v^2)^2$,
has the symmetry group $O(N)$ \emph{for the law}. The lowest-energy
solutions form the sphere $\phi_a\phi_a=v^2$; each particular solution
$\phi_a=v\,\delta_{aN}$ has the smaller isotropy group $O(N-1)$. The orbit
$O(N)/O(N-1)=S^{N-1}$ is the vacuum manifold, along which fluctuations
cost no energy: $N-1$ massless Goldstone modes. The Noether currents remain
conserved; what changes is that the solution is not invariant. This is the
distinction of Section~\ref{sec:law-solution} once more.

% ---------------------------------------------------------------------
\section{Example VI: quantum mechanics}
\label{sec:quantum}

By Wigner's theorem a symmetry of a quantum system, a bijection of rays that
preserves transition probabilities $|\langle\psi|\varphi\rangle|^2$, is
implemented by an operator $U(g)$ that is unitary or antiunitary, unique up to
a phase \cite{wigner}. The phase ambiguity means the group is represented
\emph{projectively}:
$U(g_1)U(g_2)=e^{\ii\omega(g_1,g_2)}U(g_1g_2)$. For rotations this gives
spin-$\tfrac12$ representations of the double cover $SU(2)$; for the
Galilei group the nontrivial phase is the mass $M$ of
Section~\ref{sec:galilei}, the central charge appearing in the classical
Poisson algebra.

The three notions take the following operator form.
\begin{itemize}[leftmargin=1.6em,itemsep=2pt]
\item \emph{Invariance} of the dynamics: $U(g)\,H\,U(g)^{-1}=H$, i.e.\
  $[U(g),H]=0$. Symmetry of the law is invariance of one particular operator.
\item \emph{Covariance} of observables: an observable $\mathcal O$
  transforms into $U^\dagger\mathcal OU$. A \emph{vector operator} $V_i$
  (position, momentum, angular momentum) transforms covariantly,
  \begin{equation}\label{eq:vector-op}
    U(R)^\dagger V_iU(R)=R_{ij}V_j\quad\Longleftrightarrow\quad
    [J_i,V_j]=\ii\hbar\,\epsilon_{ijk}V_k,
  \end{equation}
  whereas the scalars $H$ and $\bm J^2$ are \emph{invariant}.
\item \emph{Consequences.} If $[U,H]=0$, energy levels fall into multiplets
  that carry representations of the symmetry group (the $(2j{+}1)$-fold
  degeneracy for rotations), and matrix elements of covariant operators
  between such multiplets are constrained by the Wigner--Eckart theorem,
  yielding selection rules. Both facts follow from group theory alone,
  without solving the dynamics.
\end{itemize}
The last bullet displays the practical power of the whole subject:
covariance of operators and invariance of $H$ together determine selection
rules before any wavefunction is computed. The classical analogue of
\eqref{eq:vector-op} is $\Poi{L_i}{V_j}=\epsilon_{ijk}V_k$.

\begin{remark}[Discrete symmetries]
Parity $P$ and time reversal $T$ are improper transformations. Electrodynamics
is covariant (hence symmetric) under both if $\tilde F$, $\vB$ and $\vE$ are
assigned the appropriate parity and time-reversal signs; the weak interaction
is not symmetric under $P$, so the ``covariance'' of an equation under $P$
is a statement about the formulation only until one also asks whether the
symmetry is a property of the \emph{law} that nature obeys.
\end{remark}

% ---------------------------------------------------------------------
\section{Summary and common misunderstandings}
\label{sec:summary}

\subsection{Summary of the examples}

Table~\ref{tab:summary} places the examples in the framework of
Proposition~\ref{prop:main}.

\begin{table}[ht]
\centering
\footnotesize
\renewcommand{\arraystretch}{1.25}
\setlength{\tabcolsep}{3pt}
\begin{tabular}{@{}>{\raggedright\arraybackslash}p{2.2cm}>{\raggedright\arraybackslash}p{3.1cm}>{\raggedright\arraybackslash}p{2.7cm}>{\raggedright\arraybackslash}p{3.5cm}>{\raggedright\arraybackslash}p{3.1cm}@{}}
\toprule
System & Covariance group (of the formulation)
  & Fixed background & Symmetry $=\Stab(b_0)$ & Consequence\\
\midrule
Newtonian $N$-body & Galilei (or $\Diff$ in Newton--Cartan form)
  & $t_\mu,\ h^{\mu\nu}$ & Galilei, 10 parameters
  & $E,\bm P,\bm L,\bm G$; mass as central charge\\
Maxwell with sources & Lorentz (or $\Diff$ with $\eta$ transported)
  & $\eta_{\mu\nu}$ & Poincar\'e; conformal if $j=0$
  & $P^\mu,J^{\mu\nu}$\\
Relativistic particle & Poincar\'e, reparametrization
  & $\eta_{\mu\nu}$ & Poincar\'e; reparametrization is gauge
  & $\pi_\mu,J^{\mu\nu}$; $\pi^2-m^2c^2\weakeq0$\\
General relativity & $\Diff(M)$ & none & $\Diff(M)$, entirely gauge
  & $\nabla_\mu T^{\mu\nu}=0$; charges only for Killing $\xi$ of a solution\\
Gauge theory & local $G$ & $\eta_{\mu\nu}$
  & Poincar\'e $\times$ global $G$; local $G$ is gauge
  & charge; $\partial_\mu j^\mu\weakeq0$; Gauss law\\
Kepler problem & canonical transformations & --- 
  & $SO(3)\times\mathbb R_t$ (geometric); $SO(4)$ (dynamical, $H<0$)
  & $\vL,\ \bm A,\ H$\\
Quantum mechanics & unitary / antiunitary operators & fixed by $H$
  & $[U(g),H]=0$
  & degeneracy multiplets; selection rules\\
\bottomrule
\end{tabular}
\caption{The examples of this chapter in the language of
Proposition~\ref{prop:main}.}
\label{tab:summary}
\end{table}

\subsection{Common misunderstandings}

\begin{enumerate}[leftmargin=1.8em,itemsep=3pt,label=\textbf{\arabic*.}]
\item \textbf{``Covariant'' and ``invariant'' are not synonyms.} The fields
  $\vE,\vB$ are covariant (they mix under boosts); $E^2-B^2$ and
  $\vE\!\cdot\!\vB$ are invariant. A covariant equation $\mathcal E=0$ is
  \emph{form-invariant}, while its left-hand side is not invariant but
  transforms in a representation. The standard phrase ``the equations are
  Lorentz invariant'' is acceptable shorthand only because it is clear from
  context that it is the \emph{solution set} that is invariant.
\item \textbf{Covariance is not physics; its stabilizer is.} Any equation can
  be written covariantly (Remark~\ref{rem:kretschmann}); what matters is
  which structures must then be declared as backgrounds, and whether they are
  invariant.
\item \textbf{A symmetry of the equations need not be a symmetry of the action.}
  The rescaling $x\to\lambda x$ of a \emph{linear} equation, such as the
  harmonic oscillator, maps solutions to solutions, but changes the action by
  $\lambda^2$; it has no Noether charge. Electromagnetic duality is another case.
  The converse (variational $\Rightarrow$ equation symmetry) holds
  (Proposition~\ref{prop:variational}).
\item \textbf{A symmetry of the law is not a symmetry of the solution.}
  Elliptic orbits, the Coulomb field, the vacuum of a broken theory, and the
  Schwarzschild metric all have smaller symmetry than their laws.
\item \textbf{Gauge ``symmetry'' is a redundancy.} Constant gauge
  transformations are symmetries with charges; local gauge transformations are
  not, and ``spontaneous breaking of a local symmetry'' is a statement about
  the \emph{description} and not about physical degeneracy.
\item \textbf{Transform everything.} An active symmetry transformation must act on
  the whole closed system (fields \emph{and} sources, apparatus,
  boundary conditions). Leaving a source untransformed amounts to introducing a
  background that breaks the symmetry (Example~\ref{ex:coulomb}).
\end{enumerate}

\subsection{Notes and further reading}

Noether's original paper is \cite{noether}; a modern account of both theorems,
with their history, is \cite{ks}, and the systematic treatment of symmetries of
differential equations is \cite{olver}. The philosophical literature on
the three notions, including the debate on general covariance, is collected in
\cite{bc,earman,norton}. For the conventions used here see \cite{ll} and
\cite{jackson}.

% ---------------------------------------------------------------------
\section*{Problems}
\addcontentsline{toc}{section}{Problems}

\begin{enumerate}[leftmargin=2.2em,itemsep=4pt,label=\textbf{P\thechapter.\arabic*}]
\item Show that $\mathcal E^\nu=\partial_\mu F^{\mu\nu}-\tfrac{4\pi}{c}j^\nu$
  transforms as a four-vector under \eqref{eq:poincare}. Using
  \eqref{eq:boost-fields}, verify directly that $E^2-B^2$ and
  $\vE\!\cdot\!\vB$ are invariant.
\item For the friction law \eqref{eq:friction}, determine the subgroup of the
  Galilei group that is a symmetry when (a) $\bm w=0$, (b) $\bm w=\bm w_0$
  constant, (c) $\bm w=\bm\Omega\times\vx$ (a rigidly rotating medium).
\item Verify that $\bm G=M\bm X_{\rm cm}-\bm P\,t$ is conserved and that
  $\Poi{G_i}{P_j}=M\delta_{ij}$. How does the corresponding relation differ in
  the Poincar\'e algebra, and how does the Galilei limit produce the central
  extension?
\item Show that the harmonic oscillator equation $\ddot x+\omega^2x=0$ is
  invariant under $x\to\lambda x$ but the action is not. Show that
  $x\to x+\eps\,y(t)$, with $y$ any fixed solution, is a variational symmetry
  up to a boundary term, and compute the associated conserved quantity
  (the Wronskian).
\item Verify that the Runge--Lenz vector \eqref{eq:rl} is conserved and
  compute $\Poi{A_i}{A_j}$. Obtain the $so(4)$ algebra for $H<0$.
\item For the action \eqref{eq:particle}, derive the identity
  \eqref{eq:param-identity} and the constraint
  $\eta^{\mu\nu}\pi_\mu\pi_\nu=m^2c^2$. Fix the gauge $\lambda=t$ and
  show that the infinitesimal Lorentz boost acts \emph{nonlinearly}
  on $x^i(t)$.
\item Let $\psi\to e^{-\ii q\Lambda/\hbar c}\psi$, $A_\mu\to A_\mu+\partial_\mu\Lambda$.
  Verify \eqref{eq:covder} and $[D_\mu,D_\nu]\psi=\tfrac{\ii q}{\hbar c}F_{\mu\nu}\psi$.
  For the non-Abelian case show that $F_{\mu\nu}\to UF_{\mu\nu}U^{-1}$ and that
  $\tr F_{\mu\nu}F^{\mu\nu}$ is invariant.
\item In GR, show that if $\nabla_\mu T^{\mu\nu}=0$ and $\xi$ is a Killing vector
  then $J^\mu=T^{\mu\nu}\xi_\nu$ is conserved. Explain why a general
  asymptotically flat solution has a conserved energy only because of the
  Killing vector $\partial_t$ (or asymptotic symmetry), not because of
  general covariance.
\item Let $[U(R),H]=0$ for all rotations $R$. Prove that every energy level is
  $(2j+1)$-fold degenerate for some $j$, and that
  $\langle j'm'|V_q^{(1)}|jm\rangle$ vanishes unless
  $|j-j'|\le1\le j+j'$ and $m'=m+q$.
\end{enumerate}

% ====================== CHAPTER ENDS ==================================

\begin{thebibliography}{99}
\addcontentsline{toc}{section}{References}
\bibitem{noether} E.~Noether, \emph{Invariante Variationsprobleme},
  Nachr.\ Ges.\ Wiss.\ G\"ottingen (1918) 235--257.
\bibitem{kretschmann} E.~Kretschmann, \emph{\"Uber den physikalischen Sinn der
  Relativit\"atspostulate}, Ann.\ Phys.\ \textbf{53} (1917) 575--614.
\bibitem{einstein1918} A.~Einstein, \emph{Prinzipielles zur allgemeinen
  Relativit\"atstheorie}, Ann.\ Phys.\ \textbf{55} (1918) 241--244.
\bibitem{anderson} J.~L.~Anderson, \emph{Principles of Relativity Physics},
  Academic Press, 1967.
\bibitem{norton} J.~D.~Norton, \emph{General covariance and the foundations of
  general relativity: eight decades of dispute}, Rep.\ Prog.\ Phys.\
  \textbf{56} (1993) 791--858.
\bibitem{kn} S.~Kobayashi and K.~Nomizu, \emph{Foundations of Differential
  Geometry}, Vol.~I, Interscience, 1963.
\bibitem{sharpe} R.~W.~Sharpe, \emph{Differential Geometry: Cartan's
  Generalization of Klein's Erlangen Program}, Springer, 1997.
\bibitem{friedman} M.~Friedman, \emph{Foundations of Space-Time Theories},
  Princeton University Press, 1983.
\bibitem{kuchar} C.~J.~Isham and K.~V.~Kucha\v{r}, \emph{Representations of
  spacetime diffeomorphisms. I. Canonical parametrized field theories},
  Ann.\ Phys.\ \textbf{164} (1985) 288--315.
\bibitem{giulini} D.~Giulini, \emph{Some remarks on the notions of general
  covariance and background independence}, Lect.\ Notes Phys.\ \textbf{721}
  (2007) 105--120.
\bibitem{lovelock} D.~Lovelock, \emph{The Einstein tensor and its
  generalizations}, J.\ Math.\ Phys.\ \textbf{12} (1971) 498--501.
\bibitem{earman} J.~Earman, \emph{World Enough and Space-Time}, MIT Press, 1989.
\bibitem{cartan} \'E.~Cartan, \emph{Sur les vari\'et\'es \`a connexion affine
  et la th\'eorie de la relativit\'e g\'en\'eralis\'ee}, Ann.\ Sci.\ \'Ecole
  Norm.\ Sup.\ \textbf{40} (1923) 325--412.
\bibitem{dbkp} C.~Duval, G.~Burdet, H.~P.~K\"unzle and M.~Perrin,
  \emph{Bargmann structures and Newton--Cartan theory}, Phys.\ Rev.\ D
  \textbf{31} (1985) 1841--1853.
\bibitem{lbll} M.~Le~Bellac and J.-M.~L\'evy-Leblond, \emph{Galilean
  electromagnetism}, Nuovo Cimento B \textbf{14} (1973) 217--234.
\bibitem{deser} S.~Deser and C.~Teitelboim, \emph{Duality transformations of
  Abelian and non-Abelian gauge fields}, Phys.\ Rev.\ D \textbf{13} (1976)
  1592--1597.
\bibitem{goldstein} H.~Goldstein, C.~Poole and J.~Safko, \emph{Classical
  Mechanics}, 3rd ed., Addison-Wesley, 2002.
\bibitem{wigner} E.~P.~Wigner, \emph{Group Theory and its Application to the
  Quantum Mechanics of Atomic Spectra}, Academic Press, 1959.
\bibitem{olver} P.~J.~Olver, \emph{Applications of Lie Groups to Differential
  Equations}, Springer, 1986.
\bibitem{ks} Y.~Kosmann-Schwarzbach, \emph{The Noether Theorems}, Springer, 2011.
\bibitem{bc} K.~Brading and E.~Castellani (eds.), \emph{Symmetries in Physics:
  Philosophical Reflections}, Cambridge University Press, 2003.
\bibitem{ll} L.~D.~Landau and E.~M.~Lifshitz, \emph{The Classical Theory of
  Fields}, Pergamon Press.
\bibitem{jackson} J.~D.~Jackson, \emph{Classical Electrodynamics}, 3rd ed.,
  Wiley, 1999.
\end{thebibliography}

\end{document}
